\documentclass[12pt]{book}
\usepackage[utf8]{inputenc}
\usepackage[T1]{fontenc}
\usepackage{amsmath,amssymb,amsfonts}
\usepackage{mathrsfs}
\usepackage{amsthm}
\usepackage{tikz-cd}

% Calligraphic for categories and sheaves
\newcommand{\cat}[1]{\mathcal{#1}}
\newcommand{\sheaf}[1]{\mathcal{#1}}

% Standard math operators
\DeclareMathOperator{\Hom}{Hom}
\DeclareMathOperator{\Ext}{Ext}
\DeclareMathOperator{\Tor}{Tor}
\DeclareMathOperator{\id}{id}
\DeclareMathOperator{\dimk}{dim_k}
\DeclareMathOperator{\Rfst}{\mathbf{R}f_*}
\DeclareMathOperator{\Lfst}{\mathbf{L}f_*}
\DeclareMathOperator{\RHom}{\mathbf{R}\mathcal{H}\textit{om}}

\begin{document}

\setcounter{page}{252}
\chapter{V. DUALIZING COMPLEXES AND LOCAL DUALITY}

Introduction.

In this chapter we discuss dualizing complexes on a locally noetherian prescheme \(X\).
A dualizing complex is a complex \(R^{\bullet} \in D^{+}(X)\) such that the functor
\[
D: M^{\bullet} \to \RHom^{\bullet}\left(M^{\bullet}, R^{\bullet}\right)
\]
induces an auto-duality of the category \(D_{c}^{b}(X)\) consisting of those bounded complexes in \(D^{+}(X)\) which have coherent cohomology.
We will show that a large class of preschemes admits dualizing complexes, and that they are almost uniquely determined.

The notion of dualizing complex will allow us to write the duality theorem in a new way.

For example, let \(X\) be projective \(n\)-space over a field \(k\).
Then \(R^{\bullet}=\omega[n]\) (where \(\omega=\omega_{X/k}\) is the sheaf of relative \(n\)-differentials [III 1])
is a dualizing complex for \(X\), and \(k\) is a dualizing complex for \(k\).
We have a canonical isomorphism \(\Rfst(\omega[n]) \cong k\) [III 3.4], and hence, for any complex \(F^{\bullet} \in D_{c}^{b}(X)\), a homomorphism
\[
\theta: \Rfst \underline{D}_{X}\left(F^{\bullet}\right) \to \underline{D}_{k}\left(\Rfst\left(F^{\bullet}\right)\right),
\]

\setcounter{page}{253}
where \(\underline{D}_{X}\) and \(\underline{D}_{k}\) are the dualizing functors on \(X\) and on \(k\), respectively.
The duality theorem [III 5.1] says that \(\theta\) is an isomorphism, i.e., that the dualizing functors on \(X\) and on \(k\) commute with \(\Rfst\).

In the latter part of the chapter we discuss the local duality theorem, as in [LC or [SGA 62, expos\'e 2,4].
The proof of the duality theorem given here is different, however.

The reader will see that the existence of a dualizing complex on a locally noetherian prescheme \(X\) implies that \(X\) has finite Krull dimension.
Hence we will give a more general notion, that of a pointwise dualizing complex, to cover the case of infinite Krull dimension.

For convenience we defer the question of existence of dualizing complexes until the end of the chapter (section 10).

\end{document}